\documentclass{CSUthesis}

% AM-CV : Latex CV
% Version 1.0
% Copyright Andrew Mellor 2016 <mellor91@hotmail.co.uk>
%
% Created using a combination of these two freely available CV templates
% https://www.sharelatex.com/templates/cv-or-resume/professional-cv
% https://www.sharelatex.com/templates/cv-or-resume/kjh-vita

\ProvidesPackage{info}

\def\myauthor{John Smith}
\def\mytitle{Curriculum Vitae}

\def\mycopyright{\myauthor}
\def\mykeywords{}
\def\mybibliostyle{plain}
\def\mybibliocommand{}
\def\mysubtitle{}

\def\myaffiliation{University of Latex}
\def\mydepartment{School of CV Writing}
\def\mydepartmentemail{http://www.cv.latex.ac.uk/}
\def\myaddrone{Address One}
\def\myaddrtwo{Address Two}
\def\myaddrthree{Address Three}


\def\myemail{john.smith@email.com}
\def\myweb{johnsmith.com}
\def\myphone{00000 000000}
\def\mytwitter{@jsmith}

\def\myversion{}
\def\myrevision{}

% lipsum
\newcommand{\lipsum} {
    待编辑
}

\begin{document}
%%%%%%%%%%%%%%%%%%%%%%%%%%%%%%%%%%%%%%%%%%%%%%%%%%
% 封面
% -----------------------------------------------%
\makecoverpage

%%%%%%%%%%%%%%%%%%%%%%%%%%%%%%%%%%%%%%%%%%%%%%%%%%
% 前置部分的页眉页脚设置
% -----------------------------------------------%
\newpage
% 正文和后置部分用阿拉伯数字编连续码，前置部分用罗马数字单独编连续码（封面除外）。
% 设置封面页后的页码
\pagenumbering{Roman} % 大写罗马字母
\setcounter{page}{1} % 从1开始编号页码
% 设置页眉和页脚 
\pagestyle{fancy}
% 正文以前部分无需页眉
\fancyhf{} % 清空原有格式
\renewcommand{\headrulewidth}{0pt}
% 封面页无需页码，其他前置部分需要（按此理解扉页也是要页码的）。
\fancyhf[CF]{\thepage} % 所有（奇数和偶数）中间页脚

\ifblindreview	% 盲审不需要扉页和声明页
\else
%%%%%%%%%%%%%%%%%%%%%%%%%%%%%%%%%%%%%%%%%%%%%%%%%%
% 扉页 
% -----------------------------------------------%
\maketitlepage
\newpage

%%%%%%%%%%%%%%%%%%%%%%%%%%%%%%%%%%%%%%%%%%%%%%%%%%
% 声明页
% -----------------------------------------------%
\announcement
\newpage
\fi
%%%%%%%%%%%%%%%%%%%%%%%%%%%%%%%%%%%%%%%%%%%%%%%%%%
% 中文摘要
% -----------------------------------------------%
% 设置中文摘要
\keywordscn{中南大学；学位论文；LaTeX模板}
\categorycn{TP391}
\begin{abstractcn}
LaTeX利用设置好的模板，可以编译为格式统一的pdf。目前国内大多出版社与高校仍在使用word，word由于其强大的功能与灵活性，在新手面对形式固定的论文时，排版、编号、参考文献等简单事务反而会带来很多困难与麻烦，对于一些需要通篇修改的问题，要想达到LaTeX的效率，对word使用者来说需要具有较高的技能水平。

为了能把主要精力放在论文撰写上，许多国际期刊和高校都支持LaTeX的撰写与提交，新手不需要关心格式问题，只需要按部就班的使用少数符号标签，即可得到符合要求的文档。且在需要全篇格式修改时，更换或修改模板文件，即可直接重新编译为新的样式文档，这对于word新手使用word的感受来说是不可思议的。

本项目的目的是为了创建一个符合中南大学研究生学位论文（博士）撰写规范的TeX模板，解决学位论文撰写时格式调整的痛点。

% 图X幅，表X个，参考文献X篇（四号宋体）

\end{abstractcn}
\newpage

%%%%%%%%%%%%%%%%%%%%%%%%%%%%%%%%%%%%%%%%%%%%%%%%%%
% 英文摘要
% -----------------------------------------------%
\keywordsen{CSU;~LaTeX;~Template}
\categoryen{TP391}
\begin{abstracten}
LaTeX can be compiled into a pdf of uniform format using the set template.At present, most domestic publishers and universities still use word. Because of its powerful function and flexibility, when faced with fixed-form papers by novices, simple matters such as typesetting, numbering, and reference documents will bring many difficulties and troubles. For some problems that need to be modified throughout, to achieve the efficiency of LaTeX, it requires a high level of skill for word users.

In order to focus on the writing of papers, many international journals and universities support the writing and submission of LaTeX. Novices don't need to care about formatting issues. They only need to use a few symbolic labels step by step to get the documents that meet the requirements. And when you need to modify the entire format, you can directly recompile the template file by replacing or modifying the template file. This is incredible for the word novice to use the word.

The purpose of this project is to create a TeX template that meets the specifications of the graduate degree thesis (PhD) of Central South University, and to address the pain points of format adjustment during the dissertation writing.
\end{abstracten}
\newpage

%%%%%%%%%%%%%%%%%%%%%%%%%%%%%%%%%%%%%%%%%%%%%%%%%%
% 正文页眉页脚
% -----------------------------------------------%
\setcounter{page}{1} % 重置目录页码为小写罗马字体
\pagenumbering{roman} % 设置页码为小写罗马字体
% 设置页眉和页脚 %
\pagestyle{fancy}
\fancyhf[CF]{\thepage} % 所有（奇数和偶数）中间页脚

% 目录
% -------------------------------------------%
{
\renewcommand{\contentsname}{\hfill \heiti \zihao{3} 目\quad 录\hfill}  
	\renewcommand*{\baselinestretch}{1.0}   % 行间距
    \tableofcontents
}
\newpage
% 去掉页眉章节序号后面的“.”
\renewcommand{\sectionmark}[1]{\markright{\thesection~ #1}} 


\renewcommand{\headrulewidth}{1pt}
\fancyhf[RH]{\songti \zihao{5} \rightmark} % 设置所有（奇数和偶数）右侧页眉
\fancyhf[LH]{\songti \zihao{5} 博士学位论文} % 设置所有（奇数和偶数）左侧页眉
% 正文内容 
% --------------------------------------------%
\setcounter{page}{1} % 重置页码编号
\pagenumbering{arabic} % 设置页码编号为阿拉伯数字

% 可以使用include命令导入tex文件，从而避免过多修改本文件。

% 论文正文是主体，主体部分应从另页右页开始，每一章应另起页。一般由序号标题、文字叙述、图、表格和公式等五个部分构成。

% 重新设置正文行间距，因为前置部分设置时候行间距被改过
\renewcommand*{\baselinestretch}{1.0}   % 几倍行间距
\setlength{\baselineskip}{20pt}         % 基准行间距

% 正文
{
% 表格字号应比正文小，一般五号/10.5pt，但是暂时没法再cls里设置（不然会影响到封面等tabular环境）
% 所以目前只好在主文件里局部\AtBeginEnvironment
	\AtBeginEnvironment{tabular}{\zihao{5}}
	% Contents
\cleardoublestylepage{empty}
\thispagestyle{common}
\setcounter{page}{1}

% set as first page if inside `proposal`
\ifthenelse{\equal{\Period}{proposal}}
{
    \setcounter{page}{1}
}
{}

\ifthenelse{\equal{\Type}{thesis}}
{
    % `thesis` content
    \cleardoublepage
\chapter{文献综述}

\section{背景介绍}

\subsection{小节}

\subsection{小节}

\section{国内外研究现状}

\subsection{研究方向及进展}

\subsection{存在问题}

\section{研究展望}

    %%
% 开题报告
% modifier: 黄俊杰(huangjj27, 349373001dc@gmail.com)
% update date: 2017-05-14

% 选题目的
\objective{

}

% 思路
\methodology{

}

% 研究方法/程序/步骤
\researchProcedure{

}

% 相关支持条件
\supportment{

}

% 进度安排
\schedule{

}

% 指导老师意见
\proposalInstructions{

}


    \cleardoublepage
\chapter{外文翻译}

\section*{摘要}

\section{节标题}

    \cleardoublepage
\chapter{外文原文}

% Insert pdf version of original document
% Place the pdf under `/figure/`

\begin{center}
    \includepdf[
        pages=1-3,
        scale=0.75,
        pagecommand={\pagestyle{fancy}},
        frame, % remove this line if you don't want outer fram of pdf
    ]{rules2018.pdf}
\end{center}

}
{
    \ifthenelse{\equal{\Type}{design}}
    {
        % `design` content
        %%
% 开题报告
% modifier: 黄俊杰(huangjj27, 349373001dc@gmail.com)
% update date: 2017-05-14

% 选题目的
\objective{

}

% 思路
\methodology{

}

% 研究方法/程序/步骤
\researchProcedure{

}

% 相关支持条件
\supportment{

}

% 进度安排
\schedule{

}

% 指导老师意见
\proposalInstructions{

}


        \cleardoublepage
\chapter{外文翻译}

\section*{摘要}

\section{节标题}

        \cleardoublepage
\chapter{外文原文}

% Insert pdf version of original document
% Place the pdf under `/figure/`

\begin{center}
    \includepdf[
        pages=1-3,
        scale=0.75,
        pagecommand={\pagestyle{fancy}},
        frame, % remove this line if you don't want outer fram of pdf
    ]{rules2018.pdf}
\end{center}

        \cleardoublepage
\chapter{中期检查}

\section{项目概况}

\section{工作进展情况}

\section{问题与建议}

\section{其他}

    }
    {
        % otherwise, blank
    }
}

}

%%%%%%%%%%%%%%%%%%%%%%%%%%%%%%%%%%%%%%%%%%%%%%%%%%
% 临时标签，用于编译时追踪正文末尾
%%%%%%%%%%%%%%%%%%%%%%%%%%%%%%%%%%%%%%%%%%%%%%%%%%

%%%%%%%%%%%%%%%%%%%%%%%%%%%%%%%%%%%%%%%%%%%%%%%%%%
% 后续内容，标题三号黑体居中，章节无编号
% --------------------------------------------%

% https://www.zhihu.com/question/29413517/answer/44358389 %
% 说明如下：
% secnumdepth 这个计数器是 LaTeX 标准文档类用来控制章节编号深度的。在 article 中，这个计数器的值默认是 3，对应的章节命令是 \subsubsection。也就是说，默认情况下，article 将会对 \subsubsection 及其之上的所有章节标题进行编号，也就是 \part, \section, \subsection, \subsubsection。LaTeX 标准文档类中，最大的标题是 \part。它在 book 和 report 类中的层级是「-1」，在 article 类中的层级是「0」。这里，我们在调用 \appendix 的时候将计数器设置为 -2，因此所有的章节命令都不会编号了。不过，一般还是会保留 \part 的编号的。所以在实际使用中，将它设置为 0 就可以了。

% 在修改过程中请注意不要破环命令的完整性

\renewcommand\appendix{\setcounter{secnumdepth}{-2}}
\appendix

% 主文件有代码去掉页眉章节编号的“.”，但这会因为bug导致无编号章节显示一个错误编号，所以这里在无编号章节之前再次重定义sectionmark。
\renewcommand{\sectionmark}[1]{\markright{#1}}

% section 标题从这里往后改为三号黑体居中
\titleformat{\section}{\centering \zihao{3}\heiti}{\thesection}{1em}{}

% \section{参考文献} % bibliography会自动显示参考文献四个字
\addcontentsline{toc}{section}{参考文献} % 由于参考文献不是section，这句把参考文献加入目录
% \nocite{*} % 该命令用于显示全部参考文献，即使文中没引用
% cls文件中已经引入package，这里不需要调用 \bibliographystyle 了。
% \bibliographystyle{gbt7714-2005} 
\bibliography{thesis-references}
\newpage

\section{攻读学位期间主要研究成果} % 无章节编号

\ifblindreview
\noindent
（盲审隐去作者相关具体信息）
\fi
\subsection*{学术论文}

\ifblindreview

\noindent
第一作者：JCR 1 区 x 篇，会议 x 篇 \\{}
第二作者：JCR 1 区 x 篇，3 区 x 篇，4 区 x 篇，EI x 篇 

\noindent
投稿状态： 
IEEE Transactions on Image Processing 1篇（under review）\\{}
IEEE Transactions on Circuits and Systems for Video Technology 1 篇（Accept with Minor Revision）
\else
% 标准版本
\begin{enumerate}[label={[\arabic*]}]
\item \textbf{Daxia Guo}, Director, Someone. CSU Latex Template[J]. CSU player: 1(1):1-10. （SCI 检索，JCR 1 区）
\item Director, \textbf{Daxia Guo}, Someone, Someother. XXXXXX[J]. Transactions on Image Processing.（SCI Under Review，JCR 1 区）
\item Director, \textbf{Daxia Guo}, Someone, Someother. XXXXXX[J]. Transactions on Circuits and Systems for Video Technology.（SCI Under Review，JCR 1 区）
\end{enumerate}
\fi
\subsection*{发明专利}
\ifblindreview
\noindent
发明专利 2 项，已公开
\else
\begin{enumerate}[label={[\arabic*]}]
\item 郭大侠，XXX，XXX. 一种用Latex写中南大学学位论文的方法. 申请号：CN20190415xxxx，公开号：CNXXXXXXXXXA
\end{enumerate}
\fi

\ifblindreview
\else
\subsection*{主持和参与的科研项目}
\begin{enumerate}[label={[\arabic*]}]
\item 国家自然科学基金面上项目《XXXXXXXXXXXX》， 项目编号：XXXXXXXX，参与.
\end{enumerate}

\subsection*{个人获奖情况}
\noindent
\begin{enumerate}[label={[\arabic*]}]
\item XX金奖
\item XX奖学金
\end{enumerate}
\fi

\newpage

\ifblindreview
\else
\section{致谢} % 无章节编号
\lipsum
\newpage
\fi

\end{document}
